% SOURCE: original Chinese PDF, viewer pages 16--18.
\section{Transformation of Representations}
Selecting a suitable representation can simplify a calculation. The transformation between basis vectors of different representations is called a transformation of representation.

Consider two representations:
\begin{equation}
\begin{array}{ll}
F\text{ representation:}&\text{basis }\ket i,\quad \op F\ket i=f_i\ket i,\\
G\text{ representation:}&\text{basis }\ket m,\quad \op G\ket m=g_m\ket m.
\end{array}
\end{equation}
We have
\begin{equation}
\ket m=\sum_i\ket i\braket i m=\sum_i S_{im}\ket i.
\end{equation}
The transformation matrix $S$ is defined by
\begin{equation}
S_{im}=\braket i m.
\end{equation}
It represents the basis vectors of the $M$ representation in the $I$ representation.
\subsection*{1) Transformation of an arbitrary vector}
In the $I$ representation,
\begin{equation}
\ket a=\sum_i\ket i\braket i a.
\end{equation}
The coefficient $\braket i a$ is the component of $\ket a$ along $\ket i$, or its representation in the $I$ basis.

In the $M$ representation,
\begin{equation}
\ket a=\sum_m\ket m\braket m a.
\end{equation}
The coefficient $\braket m a$ is the component along $\ket m$, or the representation of $\ket a$ in the $M$ basis.

Using completeness,
\begin{equation}
\braket m a=\sum_i\braket m i\braket i a
=\sum_i S_{im}^*\braket i a
=\sum_i S_{mi}^\dagger\braket i a.
\end{equation}
In matrix form,
\begin{equation}
\boxed{a_G=S^\dagger a_F.}
\end{equation}
\subsection*{2) Transformation of an arbitrary operator}
In the $I$ representation,
\begin{equation}
(A_F)_{ij}=\mel i{\op A}j.
\end{equation}
In the $M$ representation,
\begin{equation}\begin{aligned}
(A_G)_{mn}=\mel m{\op A}n
&=\sum_{ij}\braket m i\mel i{\op A}j\braket j n\\
&=\sum_{ij}S_{im}^*(A_F)_{ij}S_{jn}
=\sum_{ij}(S^\dagger)_{mi}(A_F)_{ij}S_{jn}.
\end{aligned}\end{equation}
In matrix form,
\begin{equation}
\boxed{A_G=S^\dagger A_F S.}
\end{equation}
\subsection*{3) The representation-transformation matrix is unitary}
\begin{equation}
(SS^\dagger)_{ij}=\sum_m S_{im}(S^\dagger)_{mj}
=\sum_m S_{im}S_{jm}^*
=\sum_m\braket i m\braket m j
=\braket i j=\delta_{ij}.
\end{equation}
Therefore $SS^\dagger=\mathbf1$. Similarly, $S^\dagger S=\mathbf1$, and hence
\begin{equation}
\boxed{S^\dagger=S^{-1}.}
\end{equation}
A representation-transformation matrix is unitary.

Note that $S$ is not necessarily Hermitian: $S^\dagger\ne S$.
\subsection*{4) Transformation of representation does not change eigenvalues}
Suppose
\begin{equation}
A_F a_F=\lambda_F a_F.
\end{equation}
Then
\begin{equation}\begin{aligned}
A_G a_G
&=S^\dagger A_F SS^\dagger a_F
=S^\dagger A_Fa_F=S^\dagger\lambda_Fa_F\\
&=\lambda_F S^\dagger a_F=\lambda_F a_G.
\end{aligned}\end{equation}
Thus the eigenvalue in the $M$ representation is still $\lambda_F$. A transformation of representation does not change the values of physical quantities.

\textbf{Example 1:} A transformation of representation does not change commutation relations.

If $[A_I,B_I]=C_I$ in the $I$ representation, then in the $M$ representation,
\begin{equation}\begin{aligned}
[A_M,B_M]&=A_MB_M-B_MA_M\\
&=S^\dagger A_I SS^\dagger B_I S-S^\dagger B_I SS^\dagger A_I S\\
&=S^\dagger(A_IB_I-B_IA_I)S=S^\dagger C_I S=C_M.
\end{aligned}\end{equation}
Commutation relations are fundamental relations of quantum mechanics and are unchanged by the transformation of representation.

\textbf{Example 2:} A transformation of representation does not change the trace of a matrix:
\begin{equation}
\operatorname{tr}(A_G)=\operatorname{tr}(S^\dagger A_FS)
=\operatorname{tr}(SS^\dagger A_F)=\operatorname{tr}(A_F).
\end{equation}

\interactiveguide{IC-01-05}{Measurement axes and basis changes}{After the unitary-change-of-basis derivation, rotate the spin measurement axes and transform the state components. Verify that the matrix entries change while the probabilities for a fixed physical experiment remain invariant.}
